\documentclass[11pt,tikz,border=4pt]{standalone}
\usepackage[UTF8,scheme=plain,fontset=none]{ctex}
\usepackage{amsmath,amssymb}
\setmainfont{TeX Gyre Termes}
\setsansfont{TeX Gyre Heros}
\setCJKmainfont{Noto Serif SC}
\setCJKsansfont[BoldFont=SimHei]{Noto Sans SC}
\usetikzlibrary{arrows.meta,positioning,calc,fit,backgrounds,decorations.pathreplacing,shapes.geometric}

\definecolor{Ink}{HTML}{111111}
\definecolor{Gray}{HTML}{555555}
\definecolor{Light}{HTML}{F7F7F7}

\tikzset{
  wire/.style={-{Latex[length=2.0mm,width=1.25mm]},draw=Ink,line width=0.78pt},
  plain/.style={draw=Ink,line width=0.78pt},
  state/.style={-{Latex[length=1.8mm,width=1.15mm]},draw=Ink,line width=0.68pt},
  control/.style={-{Latex[length=1.7mm,width=1.05mm]},draw=Gray,densely dashed,line width=0.58pt},
  map/.style={draw=Gray,densely dotted,line width=0.58pt},
  opblock/.style={draw=Ink,fill=white,minimum width=12mm,minimum height=8mm,align=center,inner sep=2pt,line width=0.72pt},
  wideblock/.style={opblock,minimum width=27mm,minimum height=11mm},
  delay/.style={opblock,minimum width=12mm,minimum height=8mm},
  reg/.style={opblock,minimum width=11mm,minimum height=8mm},
  mux/.style={draw=Ink,fill=white,trapezium,trapezium left angle=70,trapezium right angle=110,
              minimum width=11mm,minimum height=12mm,align=center,inner sep=1.5pt,line width=0.72pt},
  gain/.style={draw=Ink,fill=white,isosceles triangle,isosceles triangle apex angle=60,
               shape border rotate=90,minimum height=10mm,minimum width=8mm,inner sep=0pt,line width=0.72pt},
  sum/.style={draw=Ink,fill=white,circle,minimum size=6.2mm,inner sep=0pt,line width=0.78pt},
  tap/.style={circle,fill=Ink,inner sep=1.25pt},
  note/.style={font=\sffamily\scriptsize,text=Gray,align=center},
  label/.style={font=\sffamily\scriptsize,text=Ink,align=center,fill=white,inner sep=1pt},
  tinylabel/.style={font=\sffamily\tiny,text=Gray,align=center,fill=white,inner sep=0.8pt},
  title/.style={font=\sffamily\bfseries\large,text=Ink},
  subtitle/.style={font=\sffamily\small,text=Gray},
  resource/.style={draw=Gray,fill=Light,rounded corners=1.2pt,minimum height=9mm,align=center,
                   inner sep=2pt,line width=0.55pt,font=\sffamily\scriptsize},
  group/.style={draw=Gray,rounded corners=1.5mm,dashed,line width=0.48pt,inner sep=3mm},
  sign/.style={font=\sffamily\bfseries\scriptsize,text=Ink,inner sep=0pt}
}

\newcommand{\sumnode}[2]{\node[sum] (#1) at #2 {$\Sigma$};}
\newcommand{\delaynode}[3]{\node[delay] (#1) at #2 {$z^{-1}$\\[-1pt]{\tiny #3}};}

\begin{document}
\begin{tikzpicture}[x=1cm,y=1cm,font=\sffamily\small,
 label/.append style={font=\small},note/.append style={font=\small},tinylabel/.append style={font=\footnotesize}]

\node[title] at (8.25,8.45) {16倍三级CIC插值器的算法级等效信流图};
\node[subtitle] at (8.25,8.02)
 {低速 $C^3\!\rightarrow\!\uparrow16\!\rightarrow\!I^3\;\equiv\;$低速 $C^2\!\rightarrow\!\mathrm{Hold}_{16}\!\rightarrow\!I^2$};

% Upper row: low-rate combs and rate change.
\node[opblock,minimum width=16mm] (x) at (0.70,5.55) {$x_8[n]$\\21 bit};
\sumnode{c1}{(3.00,5.55)}
\draw[wire] (x) -- (c1); \coordinate (xt) at (1.85,5.55); \node[tap] at (xt) {};
\delaynode{zc1}{(1.85,3.95)}{低速事件}
\draw[plain] (xt) -- (zc1.north); \draw[state] (zc1.east) -| (c1.south);
\node[sign] at ($(c1.west)+(-0.18,0.22)$) {$+$}; \node[sign] at ($(c1.south)+(-0.18,-0.23)$) {$-$};

\sumnode{c2}{(6.00,5.55)}
\draw[wire] (c1) -- (c2); \coordinate (d1t) at (4.50,5.55); \node[tap] at (d1t) {};
\node[label,above=4pt] at (4.50,5.55) {$d_1[n]$，22 bit};
\delaynode{zc2}{(4.50,3.95)}{低速事件}
\draw[plain] (d1t) -- (zc2.north); \draw[state] (zc2.east) -| (c2.south);
\node[sign] at ($(c2.west)+(-0.18,0.22)$) {$+$}; \node[sign] at ($(c2.south)+(-0.18,-0.23)$) {$-$};

\node[wideblock,minimum width=34mm,minimum height=13mm] (hold) at (9.30,5.55)
 {$\mathrm{Hold}_{16}$\\$h[16n+r]=d_2[n]$\\{\small $r=0,\ldots,15$（重复保持）}};
\draw[wire] (c2) -- (hold); \node[label,above=4pt] at (7.55,5.55) {$d_2[n]$，23 bit};
\node[opblock,minimum width=13mm] (Aout) at (12.15,5.55) {A\\$h[k]$};
\draw[wire] (hold) -- (Aout);
\node[anchor=west] at (1.30,6.55) {(a) 低速二阶差分与16倍保持};
\node[note,anchor=west] at (13.20,5.90)
 {$d_1[n]=x_8[n]-x_8[n-1]$\\$d_2[n]=d_1[n]-d_1[n-1]$};

% Lower row: high-rate integrators and output quantization.
\node[opblock,minimum width=13mm] (Ain) at (0.70,1.65) {A\\$h[k]$};
\sumnode{i1}{(3.00,1.65)}
\draw[wire] (Ain) -- (i1); \coordinate (s1t) at (4.05,1.65); \node[tap] at (s1t) {};
\delaynode{zi1}{(4.05,0.20)}{26 bit}
\draw[plain] (i1) -- (s1t); \draw[plain] (s1t) -- (zi1.north); \draw[state] (zi1.west) -| (i1.south);
\node[sign] at ($(i1.west)+(-0.18,0.22)$) {$+$}; \node[sign] at ($(i1.south)+(-0.18,-0.23)$) {$+$};

\sumnode{i2}{(7.00,1.65)}
\draw[wire] (s1t) -- (i2); \node[label,above=4pt] at (5.55,1.65) {$s_1[k]$，26 bit};
\coordinate (s2t) at (8.05,1.65); \node[tap] at (s2t) {};
\delaynode{zi2}{(8.05,0.20)}{29 bit}
\draw[plain] (i2) -- (s2t); \draw[plain] (s2t) -- (zi2.north); \draw[state] (zi2.west) -| (i2.south);
\node[sign] at ($(i2.west)+(-0.18,0.22)$) {$+$}; \node[sign] at ($(i2.south)+(-0.18,-0.23)$) {$+$};

\node[wideblock,minimum width=27mm,minimum height=12mm] (q) at (10.70,1.65)
 {舍入与饱和\\中点远离零，$>>>8$\\29$\rightarrow$20 bit};
\draw[wire] (s2t) -- (q); \node[label,above=4pt] at (8.95,1.65) {$s_2[k]$，29 bit};
\node[opblock,minimum width=18mm] (y) at (13.50,1.65) {$y_{128}[k]$\\20 bit，$128F_s$};
\draw[wire] (q) -- (y);
\node[anchor=west] at (1.30,2.80) {(b) 高速二阶积分与定点输出};
\node[note,anchor=west] at (14.60,2.10)
 {$s_1[k]=s_1[k-1]+h[k]$\\$s_2[k]=s_2[k-1]+s_1[k]$};
\node[note,anchor=west] at (14.60,0.60)
 {$z_\ell^{-1}$/$z_h^{-1}$为事件延时；\\全链为同一$128F_s$主时钟。};

\end{tikzpicture}
\end{document}
